% Amplificador Diferencial con BJT PNP
% Basado en: S03-1 Amplificador diferencial con BJT PNP
% Compilar en Overleaf o con pdflatex
% Requiere: circuitikz

\documentclass[border=10pt]{standalone}
\usepackage[utf8]{inputenc}
\usepackage{circuitikz}

\begin{document}

\begin{circuitikz}[
    american,
    scale=1.2,
    transform shape,
    every node/.style={font=\small}
]

    % === Fuente de alimentación +Vcc (arriba) ===
    \draw (0, 7) node[above]{$+V_{CC}$} to[short, *-] (0, 6.5);

    % === Resistencia de cola REE ===
    \draw (0, 6.5) to[R, l=$R_{EE}$, i<=$I_{EE}$] (0, 5);

    % === Nodo de emisores comunes ===
    \draw (0, 5) -- (-2.5, 5) -- (-2.5, 4.5);  % rama a Q1
    \draw (0, 5) -- ( 2.5, 5) -- ( 2.5, 4.5);  % rama a Q2

    % === Transistor Q1 (PNP, izquierda) ===
    % PNP: emisor arriba, colector abajo (corriente baja)
    \draw (-2.5, 4.5) node[pnp, anchor=E, xscale=-1](Q1){};
    \node[left] at (Q1.text) {$Q_1$};

    % === Transistor Q2 (PNP, derecha) ===
    \draw (2.5, 4.5) node[pnp, anchor=E](Q2){};
    \node[right] at (Q2.text) {$Q_2$};

    % === Entradas V1 y V2 (bases) ===
    \draw (Q1.B) -- ++(-1, 0) node[left]{$V_1$};
    \draw (Q2.B) -- ++( 1, 0) node[right]{$V_2$};

    % === Resistencias de colector RC1 y RC2 ===
    \draw (Q1.C) to[R, l_=$R_{C1}$] ++(0, -2.5) node[ground](gnd1){};
    \draw (Q2.C) to[R, l=$R_{C2}$]  ++(0, -2.5) node[ground](gnd2){};

    % === Salidas Vo1 y Vo2 ===
    \draw (Q1.C) to[short, *-] ++(-1.2, 0) node[left]{$V_{o1}$};
    \draw (Q2.C) to[short, *-] ++( 1.2, 0) node[right]{$V_{o2}$};

    % === Etiquetas de corriente ===
    \draw[->, >=stealth] (-2.5, 5) -- (-2.5, 4.7) node[midway, left]{$I_{E1}$};
    \draw[->, >=stealth] ( 2.5, 5) -- ( 2.5, 4.7) node[midway, right]{$I_{E2}$};

    % === Anotaciones ===
    \node[align=center, font=\footnotesize, text=gray] at (0, 3)
        {$I_{EE} = I_{E1} + I_{E2}$};

\end{circuitikz}

\end{document}
